\section{SGLang 的位置：AI 五层架构里的中间层}

赵晨阳开场先介绍自己：清华计算机本科，2024 年毕业，在 UCLA 读了十五个月博士，随后在 2025 年 11 月左右开始创业，加入推理引擎方向的公司。主持人随即把问题引向 SGLang：它如何从开源组织走向公司化，赵晨阳在其中有什么体会。

他先解释什么是推理引擎。普通用户接触 AI，看到的是 ChatGPT、Claude、AI PPT、Sora、AI Coding 等前端产品；但用户请求发到后端之后，服务器如何调度模型、如何管理 KV cache、如何量化、如何处理不同模型形态、如何提高吞吐和降低延迟，这才是推理引擎要解决的问题。Hugging Face 的 \texttt{model.generate} 可以帮助人理解最朴素的生成过程，但真正服务线上请求，需要大量系统层优化。

赵晨阳借用一个五层架构来定位推理：最上层是 Agent 和应用；下面是具体模型，如 Claude、GPT、开源模型；再往下是推理引擎，如 OpenAI internal engine、TensorRT-LLM、SGLang；再往下是硬件，再往下是半导体、电力等基础设施。越往上变化越快，越往下周期越慢；推理引擎恰好处在中间，是模型厂商和硬件厂商之间的关键承接层。

\begin{knowledgebox}{推理引擎为什么重要}
\begin{tabularx}{\textwidth}{>{\raggedright\arraybackslash}p{0.24\textwidth} Y}
\toprule
层次 & 推理引擎要处理的问题 \\
\midrule
请求调度 & 不同用户、不同上下文长度、不同 batch 如何高效排队。 \\
KV Cache & 长上下文和多轮对话中，缓存如何复用、搬运和释放。 \\
模型形态 & 自回归、扩散、flow matching 等不同推理过程有不同系统需求。 \\
成本结构 & 延迟、吞吐、显存、硬件利用率共同决定产品能否规模化。 \\
生态适配 & 上层模型变化快，下层硬件变化慢，中间层必须持续吸收两边变化。 \\
\bottomrule
\end{tabularx}
\end{knowledgebox}

这一段最值得保留的，不是 SGLang 的具体技术细节，而是赵晨阳对“中间层”的判断。AI 时代越往上层越热闹，越容易被产品叙事占满；但真正决定应用能否规模化的，往往是中间层的成本、吞吐和工程可靠性。推理引擎不像模型发布那样容易制造新闻，也不像应用那样直接面对用户，却会成为整个 AI 经济能否跑通的重要地基。

\subsection{本章小结}

SGLang 所在的推理引擎层，是 AI 应用、模型和硬件之间的中间层。它不直接回答“模型有多聪明”，而回答“模型服务如何以可承受成本稳定运行”。赵晨阳的经历之所以适合讨论退学潮，是因为他参与的不是抽象论文项目，而是一个处在产业成本结构核心位置的开源系统。

